\documentclass[11pt]{article}
\usepackage[margin=1in]{geometry}
\usepackage{enumitem}
\usepackage{xcolor}
\usepackage{titlesec}
\usepackage{hyperref}

\definecolor{finteal}{RGB}{45,110,90}
\titleformat{\section}{\large\bfseries\color{finteal}}{}{0em}{}
\titlespacing{\section}{0pt}{1.2em}{0.4em}
\setlist[itemize]{leftmargin=1.4em, itemsep=2pt, topsep=2pt}
\hypersetup{colorlinks=true, urlcolor=finteal}

\title{\vspace{-1.5em}\textbf{FinLingo}\\[-0.2em]{\large\normalfont a financial-literacy simulator}}
\author{}
\date{}

\begin{document}
\maketitle
\vspace{-2.5em}

\section*{Inspiration}
Most young women start their careers having never been taught how money actually
works---budgeting, investing, debt, and negotiating a first salary. The tools that
exist tend to fall into two camps: financial-literacy apps that feel like dry quizzes,
and games that are fun but teach nothing real. We wanted to fuse them into something
genuinely addictive. FinLingo is a cozy, pixel-art life simulator built for
early-career women to \emph{build money confidence}---where every tap is a real
financial decision and your net worth is the score. Think Duolingo's bite-sized
loop meets an idle life-sim, set in a dorm room where learning literally pays.

\section*{What it does}
FinLingo is a native iOS game. You begin by creating a profile---name, age, income,
monthly spending, existing debt, and goals---and \emph{every calculation in the game is
personalized to it}. From your dorm room, two laptops drive the experience:
\begin{itemize}
  \item \textbf{Learn (left laptop).} Bite-sized lessons on budgeting (50/30/20),
  emergency funds, compound interest, credit, high-yield savings, salary negotiation,
  and 529 college savings. Each lesson mixes three interactive formats---multiple
  choice, drag-to-sort, and tap-to-match---and every correct answer pays real in-game
  cash (once, so it can't be farmed). A ``Practice this'' link jumps straight to the
  matching hands-on tool.
  \item \textbf{Simulate (right laptop).} A \emph{running, month-by-month life
  simulation}: net worth sits on the HUD, the clock ticks a month every few seconds,
  and a live sparkline tracks your trajectory. A real cash-flow model moves money each
  month---income in, an investment allocation that compounds, spending out, and debt
  that accrues interest. \textbf{Life-event scenarios} interrupt with real decisions
  (your car breaks down, a surprise bonus lands, a medical bill, a 401(k) match) whose
  choices bend the curve. Hands-on tools include a trading sandbox on \emph{real market
  history}, 401(k) and 529 projections, a debt-payoff planner, an emergency-fund
  calculator, and a ``Future You'' projection.
\end{itemize}
Along the way you can adopt a pet cat that follows you around and deck out your dorm---so
the reward loop of an idle game reinforces the financial one.

\section*{How we built it}
\begin{itemize}
  \item \textbf{Native iOS in Swift.} SwiftUI for the interface, SpriteKit for the
  explorable dorm world, and Combine to drive the simulation clock.
  \item \textbf{100\% code-drawn pixel art.} There are no image assets---the player
  sprite, the pet cat, the dorm, and the title-screen ``Freshman Year'' crest are all
  generated procedurally with Core Graphics and SwiftUI \texttt{Canvas}, so everything
  scales crisply with nearest-neighbor rendering.
  \item \textbf{The simulation engine.} A Combine timer advances one month at a time,
  applying a cash-flow model: income $\rightarrow$ investment allocation (compounding at
  a 7\% annual return) $\rightarrow$ spending $\rightarrow$ carried debt (compounding at
  $\sim$18\% APR), with any shortfall rolling into debt.
  \item \textbf{Real financial data.} The trading sandbox replays actual market history
  (e.g.\ AAPL, NVDA) sourced via Alpha Vantage.
  \item \textbf{Durable state.} Progress persists as Codable JSON with NaN-safe encoding
  and a corrupt-save quarantine so a bad file can never silently wipe a player.
  \item \textbf{Synthesized audio} for the UI clicks and a calm ambient loop---no
  licensed assets.
  \item \textbf{AI-assisted QA.} We ran repeated systematic-debugging sweeps with
  parallel review agents auditing the money math, persistence, UI, and scene layers.
\end{itemize}

\section*{Challenges we ran into}
\begin{itemize}
  \item Turning the Simulator from a set of static calculators into a genuine
  \emph{running} simulation---net-worth HUD, ticking clock, pause/speed controls, a live
  curve, and steerable life events---without it feeling like a spreadsheet.
  \item Drawing rich pixel art entirely in code: every sprite and the title diorama is
  hand-plotted rectangles, coins, and light cones.
  \item Data-loss bugs: manual \texttt{Codable} conformance silently wiped saves when a
  field was missing from the encoder/decoder, and unbounded compounding could push a
  \texttt{NaN}/\texttt{Inf} into the persisted net-worth history.
  \item Sequencing a clean first run (title $\rightarrow$ profile $\rightarrow$ tutorial
  $\rightarrow$ first scenario) and \emph{freezing the world} behind modals and the coach
  tour so time wouldn't advance out of sight.
  \item Financial correctness: money conservation, compounding order, debt amortization,
  and divide-by-zero guards across every calculator.
\end{itemize}

\section*{Accomplishments that we're proud of}
\begin{itemize}
  \item A genuinely running, month-by-month life simulation with net worth as the score
  and life events that change the outcome.
  \item Every pixel of art is generated in code---zero external image files.
  \item Real stock-market data powering a risk-free trading sandbox.
  \item Six-plus lessons across three interactive formats with abuse-proof, once-ever
  scoring, each linked to a hands-on practice tool.
  \item A cohesive, presentable first-run flow and a polished, game-style home screen.
  \item Rigor: parallel-agent debugging verified Codable symmetry, financial math, and
  crash safety end to end.
\end{itemize}

\section*{What we learned}
\begin{itemize}
  \item How to integrate SwiftUI and SpriteKit and drive a simulation loop with Combine.
  \item Real financial modeling---annuity future value, compound interest, the Rule of
  72, and debt amortization---translated into code a player can \emph{feel}.
  \item The sharp edges of hand-written \texttt{Codable} conformance and floating-point
  safety in persisted data.
  \item Procedural pixel-art techniques for characters, scenes, and UI.
  \item Designing for a specific audience: clarity, encouragement, and relevance beat
  jargon every time.
\end{itemize}

\section*{What's next for FinLingo}
\begin{itemize}
  \item \textbf{More levels beyond Freshman Year}---sophomore year, a first job and
  salary negotiation, moving out, buying a home, and long-term retirement.
  \item \textbf{A room-decoration economy}: spend earned cash on posters, plants, and
  lights, plus richer pet interactions, so the idle loop deepens.
  \item \textbf{Community}: leaderboards, shareable net-worth cards, and ``\% of people
  your age got this right'' insights.
  \item \textbf{More content and personalization}: additional life events and lessons,
  plus AI-generated lesson plans tailored to each player's goals.
  \item \textbf{New modes}: a daily money puzzle and a card-based ``Money Run'' roguelike
  for replayable, addictive learning.
  \item \textbf{Backend and cloud save} to power real social features across devices.
\end{itemize}

\end{document}
